\scanpage{156}
\mspage{74}
12. Digression sur les déploiements des diviseurs à croisements normaux (suite). \deleted{\unsure{Des propriétés}} \inserted{Topos} canoniques associés à la stratification, \deleted{des \unsure{morphismes} de groupoïdes}, leur dévissage en topos élémentaires.

\deleted{Je me dis} Même dans le cas, en apparence anodin, où $X=D_0$ est une courbe alg. \inserted{complexe} lisse (\deleted{propre} complète p.ex), et où $\Theta$ se réduit à un diviseur lisse, le grand formalisme développé dans la section précédente est incomplet. En fait, le ``th. de recollement'' est alors une pure tautologie, \deleted{\unsure{sur le}} \inserted{revenant} \deleted{voisinage} \inserted{\unsure{consistant}} à dire que le voisinage tubulaire de $\Theta$ dans $X$, privé de $\Theta$, est la somme des \deleted{voisinages} disques épointés correspondant aux pts de $\Theta$! Il y a cependant un énoncé bien intéressant, qui exprime \deleted{$\unsure{\Pi}$} $\Pi_1(X)$ (et non $\Pi_1\underbrace{X^*}_{X\setminus\Theta}\deleted{\unsure{\Theta}}$) comme somme amalgamée de
\diagram{\begin{tikzcd}[ampersand replacement=\&,column sep=small]
\&\Pi_1(V^*_{\Theta,X\deleted{\unsure{\Theta}}})\arrow[dl]\arrow[dr]\&\\
\Pi_1X^*\&\&\Pi_1V_{\Theta,X}
\end{tikzcd}}
\margin{groupoïde \uline{discret} défini par $\Theta$ comme ens. de sommets}
i.e. de
\[
(254)
\]
\diagram{\begin{tikzcd}[ampersand replacement=\&,column sep=small]
\&\widetilde\Pi_{1,0}\arrow[dl]\arrow[dr]\&\\
\widetilde\Pi_0\&\&\deleted{\unsure{\Pi_{1,0}}\text{ discret}}
\end{tikzcd}}
\margin{--}
\margin{155}
